% Source PDF page 4 - reconstructed table
\begin{frame}{Packets Processing By The Firewall}
\centering\fontsize{9.2}{11.1}\selectfont
\renewcommand{\arraystretch}{1.06}
\begin{tabular}{|Y{1.5cm}|Y{2.2cm}|Y{2.7cm}|Y{5.9cm}|}\hline
\rowcolor{MethodColor!80}\color{white}\bfseries Queue Type&\color{white}\bfseries Queue Function&\color{white}\bfseries Chain in Queue&\color{white}\bfseries Chain Function\\\hline
\textcolor{CommandGreen}{\bfseries Nat}&Network Address Translation&\textcolor{WarningColor}{\bfseries PREROUTING}&Address translation occurs before routing. Used with NAT of the destination IP address, also known as destination NAT or DNAT.\\\cline{3-4}
&&\textcolor{WarningColor}{\bfseries POSTROUTING}&Address translation occurs after routing. This implies that there was no need to modify the destination IP address of the packet as in pre-routing. Used with NAT of the source IP address using either one-to-one or many-to-one NAT. This is known as source NAT, or SNAT.\\\cline{3-4}
&&\textcolor{WarningColor}{\bfseries OUTPUT}&Network address translation for packets generated by the firewall. (Rarely used in SOHO environments)\\\hline
\end{tabular}
\end{frame}
